\section{Summable physical coercivity and critical singular closure}
\label{sec:ss-coercivity}
The maximum estimate in Theorem~\ref{thm:sc-quarter} is inherited.
Summing its fourth power would introduce a factor $m+1$. We instead
sum the positive gap before estimating the transverse work. This gives
one budget for the complete record, including arbitrarily small incidences.
Throughout, $R\in[9/20,47/100]$, $m\ge1$, and the notation
$Q,O,n,v,\tau,Y,\ell,\delta,e,L,D,E$ is exactly that of
Section~\ref{sec:sc-sublevel}. Set
\begin{equation}\label{eq:ss-constants}
 r=47/100,\quad N=94,\quad a=1/416640,\quad
 T=187/25,\quad \tau_-=3/50.
\end{equation}
No probability, roof window or return event is changed by the following
comparison.

\begin{definition}[Closure of a fixed physical word]\label{def:ss-closure}
A closed physical reference is a coordinatewise limit of actual regular
records with the same finite ordered centers $(O_0,\ldots,O_m)$, with
radii converging inside the original closed radius interval. The record
length and word are fixed while taking this closure. At a limiting
contact there may be zero incidence or an unlisted disk tangency. We
regard the limit as a polygonal reference, not as a definition of the
billiard map at that singularity, and do not assign it return labels by
continuity.
\end{definition}

\begin{lemma}[Positive gap on the closed physical reference]
\label{lem:ss-closed-gap}
Actual regular records and closed physical references satisfy the
reflection and block bounds
\begin{equation}\label{eq:ss-closed-reflection}
 v_i-v_{i-1}=2c_i n_i,\quad 0\le c_i\le1,\qquad
 \sum_{i=b+1}^{b+93}c_i\ge1/4480\quad(b+94\le m).
\end{equation}
For every $Y\in\mathcal K$, not only a minimizing polygon, put
\[
 G=A(Y)-D,\qquad
 h_i=d_i\cdot(Y_i-O_i)+R|d_i|,\qquad
 E_0=R(1-c_0),\quad E_m=R(1-c_m).
\]
Then $G\ge0$, $E=E_0+E_m=L-D$, and
\begin{align}
 G&=\tfrac12\sum_{i=0}^{m-1}\ell_i|e_i|^2+\sum_{i=0}^m h_i,
 &h_i&\ge0,\label{eq:ss-exact-gap}\\
 \sum_{i=1}^{m-1}c_i|\delta_i|^2&\le R\sum_{i=1}^{m-1}h_i,
 &\sum_i\ell_i|e_i|^2&\le2G,\label{eq:ss-energy}\\
 |\delta_0|^2&\le4R(h_0+E_0),
 &|\delta_m|^2&\le4R(h_m+E_m).\label{eq:ss-end-energy}
\end{align}
Also $|\delta_i|\le2r$ and $\tau_-\le\ell_i\le T$.
\end{lemma}
\begin{proof}
For regular records, the exact gap is the telescoping identity
\eqref{eq:sc-positive-gap}; it holds on the entire disk product.
The disk identity \eqref{eq:sc-disk-curvature} gives each interior
contact inequality, including when $c_i=0$. For the first endpoint,
use $Z_0=O_0+Rv_0$. Then
$|Y_0-Z_0|^2\le2Rh_0$ and $|Q_0-Z_0|^2=2RE_0$.
The inequality $|u+v|^2\le2|u|^2+2|v|^2$ gives
\eqref{eq:ss-end-energy}; use $Z_m=O_m-Rv_{m-1}$ at the final
endpoint. Since consecutive disks are separated, every edge is
nonzero, even on the whole comparison product. The upper bound is
$\ell_i\le\tau_i+|\delta_i|+|\delta_{i+1}|\le28/5+4r=T$.
For a closed reference, positions, unit edge directions, normals,
incidences and lengths converge: their denominators are bounded below
by $R_-$ or $\tau_-$. The regular reflection equations, horizon bound
and each fixed finite block inequality therefore pass to the limit.
The same algebra then applies directly at the limit. This passage does
not continue a flow across a grazing collision.
\end{proof}

\begin{theorem}[Complete-record quartic coercivity]
\label{thm:ss-quartic}
For every actual regular record or closed physical reference, every
$m\ge1$ and every $Y\in\mathcal K$,
\begin{equation}\label{eq:ss-global-gap}
 \sum_{i=0}^m|Y_i-Q_i|^4\le C_G G+C_E E,
\end{equation}
where the following constants are independent of the word, its length,
the original radius, and all incidences:
\begin{align}
 C_G&=\frac{4r^3}{a}+16r^3+
 8r^2\left(r+\frac{3Nr}{a}+6N^2T+12N^2r\right)
       =\frac{12318659063517}{125000},\label{eq:ss-CG}\\
 C_E&=16r^3+96Nr^3=\frac{11731999}{12500}.
 \label{eq:ss-CE}
\end{align}
In particular, with $C=10^8$, the entire inherited sublevel obeys
\begin{equation}\label{eq:ss-sublevel-sum}
 \sup_{Y\in\mathcal S(Q)}\sum_{i=0}^m|Y_i-Q_i|^4\le CE.
\end{equation}
The supremum is over the continuous compact sublevel, not a sample of
minima. This is a summability estimate, rather than the $m$-multiple of
a one-contact bound.
\end{theorem}
\begin{proof}
Call an interior index an anchor when $c_i\ge a$. Let $\mathcal A$
be the anchor set and $\mathcal I$ the other interior indices. For each
$i\in\mathcal I$, let $j(i)$ be the last anchor strictly before $i$,
or $0$ when there is none. Equation~\eqref{eq:ss-closed-reflection}
implies $i-j(i)<N$: a complete $N$-flight block starting at $j(i)$
would otherwise have an interior anchor before $i$. Thus each anchor
or initial endpoint is used as a start at most $N$ times. Each flight
or interior contact lies in at most $N$ of the intervals
$[j(i),i]$. These are deterministic overlap bounds on one record;
there is no multiplication by the number of its words.

Put $w_l=v_l^\perp$ and $b_i=w_{i-1}\cdot\delta_i$.
The exact transverse telescope \eqref{eq:sc-transverse}, which also
holds in the closure, and the inequality
$(x+y+z)^2\le3(x^2+y^2+z^2)$ give
\begin{align*}
 |b_i|^2\le{}&3|\delta_{j(i)}|^2
  +3NT\sum_{l=j(i)}^{i-1}\ell_l|e_l|^2
  +12N\sum_{l=j(i)+1}^{i-1}c_l|\delta_l|^2.
\end{align*}
Indeed the squared flight term is at most
$(\sum\ell_l)(\sum\ell_l|e_l|^2)$; the squared turning term is at
most $4(\sum c_l)(\sum c_l|\delta_l|^2)$, since
$|w_l-w_{l-1}|=2c_l$. All terms here are physical reflection
quantities of the reference, not of $Y$.

The starting-point budget is shared, not paid anew for every $i$.
Disjointness of the anchor and initial contact terms and $4r\le r/a$
imply
\begin{align*}
 \sum_{j\in\mathcal A}|\delta_j|^2+|\delta_0|^2
 &\le \frac r a\sum_{j\in\mathcal A}h_j+4r h_0+4r E_0
 \le \frac r a G+4r E.
\end{align*}
Summing the preceding transverse bound and using the overlap at most
$N$ and \eqref{eq:ss-energy} now proves
\begin{equation}\label{eq:ss-transverse-sum}
 \sum_{i\in\mathcal I}|b_i|^2
 \le\left(\frac{3Nr}{a}+6N^2T+12N^2r\right)G+12Nr E.
\end{equation}
This is where a complete-record budget replaces the unbounded number
of local estimates.

For $i\in\mathcal I$, disk curvature and
$n_i\cdot v_{i-1}=-c_i$ give
\[
 |\delta_i|^2\le-2R n_i\cdot\delta_i
 \le2r(c_i|\delta_i|+|b_i|),\qquad
 |\delta_i|^4\le8r^2(c_i^2|\delta_i|^2+|b_i|^2).
\]
Use $c_i^2\le c_i$, sum, and apply \eqref{eq:ss-energy} and
\eqref{eq:ss-transverse-sum}. The nonanchor contribution is at most
\[
 8r^2\left(r+\frac{3Nr}{a}+6N^2T+12N^2r\right)G+96Nr^3E.
\]
For anchors, $|\delta_i|^4\le4r^2|\delta_i|^2$ and the contact
budget give a total at most $4r^3G/a$. The two endpoints together
contribute at most $16r^3(G+E)$ by \eqref{eq:ss-end-energy}.
Adding these three disjoint classes proves \eqref{eq:ss-global-gap}.
It also covers $m=1$, empty anchor sets and both incomplete end blocks.
The exact rational evaluations \eqref{eq:ss-CG}--\eqref{eq:ss-CE}
satisfy
\[
 C_G+C_E=\frac{12318776383507}{125000}<10^8,
 \qquad 10^8-C_G-C_E=\frac{181223616493}{125000}>0.
\]
On $\mathcal S(Q)$, $0\le G=A(Y)-D\le L-D=E$, proving the
claimed supremum estimate. No exchange of a collision limit with that
supremum has been made.
\end{proof}

\begin{corollary}[Coercivity at a critical singular limit]
\label{cor:ss-critical-closure}
If the closed physical reference has normal endpoints, so $E=0$,
then for every $Y\in\mathcal K$
\begin{equation}\label{eq:ss-critical-coercivity}
 A(Y)\ge L+C_G^{-1}\sum_{i=0}^m|Y_i-Q_i|^4.
\end{equation}
Thus $Q$ is the unique disk-product minimum, even if some internal
incidences vanish. An approximate minimum of excess at most $\eta$
has $\sum_i|Y_i-Q_i|^4\le C_G\eta$.
\end{corollary}
\begin{proof}
Here $D=L$, so $G=A(Y)-L$. Apply
\eqref{eq:ss-global-gap}. Nonnegativity and equality force every
coordinate to agree with $Q$; the approximate assertion follows from
$G\le\eta$. In particular this proof does not divide by a vanishing
incidence, unlike the regular critical contact argument alone.
\end{proof}

\begin{corollary}[Unique fixed-word compactification]
\label{cor:ss-compactification}
Fix $m$ and its ordered centers. Any sequence of actual regular
normal-endpoint records of that word with $R_j\to R_*$ in the
original interval converges to the same closed physical reference.
Its roof equals $\ell_w(R_*)$. This asserts uniqueness of the
limiting polygon, not physical continuation, section-label continuity
or a coarea density at the singular limit.
\end{corollary}
\begin{proof}
The finite product of bounded disks is compact. Every convergent
subsequence is a closed physical reference with normal endpoints.
By Corollary~\ref{cor:ss-critical-closure}, each such limit is the
unique minimizer of the same functional at $R_*$. All subsequential
limits therefore agree, which proves convergence of the entire
sequence and the roof identity by continuity of Euclidean length.
The conclusion does not require the set of admissible radii to be
connected. No orbit is inserted at a radius where none is physical.
\end{proof}

\begin{remark}[What the singular statement does not remove]
The limiting word may have an unlisted tangency and a singular contact
can carry paired boundary data not seen by bulk collision measure.
Uniqueness of its minimizing polygon does not bound the inverse roof
Jacobian, continue the return map, or show that its normalized paired
source vanishes. All such trace terms in the inherited source remain
present. The constants above are intentionally nonoptimal; their
independence of the collision count and incidence is the conclusion.
\end{remark}
